\subsection{Jacobi transformation}
The basic Jacobi rotation $P_{pq}$ is a matrix of the form: \begin{equation}
  \mathbf{P}_{pq} = 
\begin{bmatrix}
1 & & & & & & \\
 & \dots & & & & & \\
 & & c & \dots & s & & \\
 & & \vdots & 1 & \vdots & & \\
 & & -s & \dots & c & & \\
 & & & & & \dots & \\
 & & & & & & 1
\end{bmatrix}  
\end{equation}
In rows and columns $p$ and $p$ the elements c are the only diagonal elements which are not unify. The values $c$ and $s$ are the cosine and sine of rotation angle $\phi$ so that the pythagorean theorem is fufilled($c^2 + s^2 = 1$).\\
The transformation updates the matrix $\mathbf{A}$ via an orthogonal similarity transformation:
\begin{equation}
\mathbf{A}' = \mathbf{P}_{pq}^T \mathbf{A} \mathbf{P}_{pq}
\end{equation}
With this change only $P_{pq}^T$ changes only the rows and the factor $P_{pq}$ only changes the columns. 
The idea of the Jacobi method is to try to zero the off-diagonal elements by a
series of plane rotations(\cite{explanationJacobi}). By evaluating the explicit multiplication formulas, this constraint yields the relationship:
\begin{equation}
\theta \equiv \cot 2\phi = \frac{a_{qq} - a_{pp}}{2a_{pq}}
\end{equation}
To ensure numerical stability, the algorithm solves for the smaller root of the corresponding quadratic equation $t^2 + 2t\theta - 1 = 0$ where $t \equiv s/c$. The smaller root of this equation corresponds toa rotation angle less than $\pi/4$
in magnitude; thi schoice at each stage gives the most stable reduction(\cite{explanationJacobi}).:
\begin{equation}
t = \frac{\text{sgn}(\theta)}{|\theta| + \sqrt{\theta^2 + 1}}
\end{equation}
Once the $t$ is computed the rotation matrix parameters are computed by: 
\begin{equation}
c = 1 / \sqrt{t^2+1}
\end{equation}
An important restriction is that if $\theta$ is chosen too large ar that its square $\theta^2$ whould overflow the $t$ is set to $1/2\theta$.
